\section{Introduction}

\subsection{EUV Light Source Design Criteria}

The application-level wavelength, bandwidth, flux, stability, polarization,
repetition-rate, and availability requirements for EUVICS are not yet approved.
They are tracked as candidate requirements SYS-EUV-001 through SYS-EUV-008 in
\texttt{tables/requirements\_matrix.csv}. The 13.5~nm case is retained as a
candidate design target rather than an application requirement. NIST maintains
calibration capability in the extreme ultraviolet (EUV), illustrating the need
for traceable radiometry when source performance is eventually
verified~\cite{NISTEUVDetectorCalibration}.

\subsection{Accelerator-Based EUV Light Sources}

In a Compton light source, the predicted photon distribution depends on the
electron-beam, laser-beam, interaction-geometry, and collimation inputs. Analytical
and Monte Carlo treatments of these dependencies are described by Sun and
Wu~\cite{SunWu2011ComptonSource}. This source supports the modeling dependency;
it does not establish EUVICS brightness, bandwidth, compactness, or feasibility.

\subsection{ICS EUV Light Source Design Concept and Performance Goal}

The candidate EUVICS concept brings an electron bunch and an incident laser pulse
to an interaction region and collects scattered radiation within a defined
acceptance. The nominal roadmap inputs are a 3~MeV electron \emph{kinetic} energy
and an 800~nm incident-laser wavelength. Output wavelength, flux, bandwidth,
brilliance, stability, efficiency, and acceptance remain TBD pending an approved
configuration and versioned pyEUVICS calculation. The 6.7~nm and 13.5~nm values
are candidate design targets, not predictions or measurements.

\blankline
\blankline
